\section{Integral and finite morphisms}
\label{section-integral}

\noindent
Recall that a ring map $R \to A$ is said to be {\it integral}
if every element of $A$ satisfies a monic equation with
coefficients in $R$. Recall that a ring map $R \to A$ is
said to be {\it finite} if $A$ is finite as an $R$-module.
See Algebra, Definition \ref{algebra-definition-integral-ring-map}.

\begin{definition}
\label{definition-integral}
Let $f : X \to S$ be a morphism of schemes.
\begin{enumerate}
\item We say that $f$ is {\it integral} if $f$ is affine
and if for every affine open $\Spec(R) = V \subset S$
with inverse image $\Spec(A) = f^{-1}(V) \subset X$
the associated ring map $R \to A$ is integral.
\item We say that $f$ is {\it finite} if $f$ is affine
and if for every affine open $\Spec(R) = V \subset S$
with inverse image $\Spec(A) = f^{-1}(V) \subset X$
the associated ring map $R \to A$ is finite.
\end{enumerate}
\end{definition}

\noindent
It is clear that integral/finite morphisms are separated and
quasi-compact. It is also clear that a finite morphism is a
morphism of finite type. Most of the lemmas in this section are
completely standard.
But note the fun Lemma \ref{lemma-integral-universally-closed}
at the end of the section.

\begin{lemma}
\label{lemma-integral-local}
Let $f : X \to S$ be a morphism of schemes.
The following are equivalent:
\begin{enumerate}
\item The morphism $f$ is integral.
\item There exists an affine open covering $S = \bigcup U_i$ such that
each $f^{-1}(U_i)$ is affine and
$\mathcal{O}_S(U_i) \to \mathcal{O}_X(f^{-1}(U_i))$ is integral.
\item There exists an open covering $S = \bigcup U_i$
such that each $f^{-1}(U_i) \to U_i$ is integral.
\end{enumerate}
Moreover, if $f$ is integral then for every open subscheme
$U \subset S$ the morphism $f : f^{-1}(U) \to U$ is integral.
\end{lemma}

\begin{proof}
See Algebra, Lemma \ref{algebra-lemma-integral-local}.
Some details omitted.
\end{proof}

\begin{lemma}
\label{lemma-finite-local}
Let $f : X \to S$ be a morphism of schemes.
The following are equivalent:
\begin{enumerate}
\item The morphism $f$ is finite.
\item There exists an affine open covering $S = \bigcup U_i$ such that
each $f^{-1}(U_i)$ is affine and
$\mathcal{O}_S(U_i) \to \mathcal{O}_X(f^{-1}(U_i))$ is finite.
\item There exists an open covering $S = \bigcup U_i$
such that each $f^{-1}(U_i) \to U_i$ is finite.
\end{enumerate}
Moreover, if $f$ is finite then for every open subscheme
$U \subset S$ the morphism $f : f^{-1}(U) \to U$ is finite.
\end{lemma}

\begin{proof}
See Algebra, Lemma \ref{algebra-lemma-integral-local}.
Some details omitted.
\end{proof}

\begin{lemma}
\label{lemma-finite-integral}
A finite morphism is integral.
An integral morphism which is locally of finite type is finite.
\end{lemma}

\begin{proof}
See Algebra, Lemma \ref{algebra-lemma-finite-is-integral}
and Lemma \ref{algebra-lemma-characterize-finite-in-terms-of-integral}.
\end{proof}

\begin{lemma}
\label{lemma-composition-finite}
A composition of finite morphisms is finite.
Same is true for integral morphisms.
\end{lemma}

\begin{proof}
See Algebra, Lemmas \ref{algebra-lemma-finite-transitive}
and \ref{algebra-lemma-integral-transitive}.
\end{proof}

\begin{lemma}
\label{lemma-base-change-finite}
A base change of a finite morphism is finite.
Same is true for integral morphisms.
\end{lemma}

\begin{proof}
See Algebra, Lemma \ref{algebra-lemma-base-change-integral}.
\end{proof}

\begin{lemma}
\label{lemma-integral-universally-closed}
\begin{slogan}
integral $=$ affine $+$ universally closed
\end{slogan}
Let $f : X \to S$ be a morphism of schemes.
The following are equivalent
\begin{enumerate}
\item $f$ is integral, and
\item $f$ is affine and universally closed.
\end{enumerate}
\end{lemma}

\begin{proof}
Assume (1). An integral morphism is affine by definition.
A base change of an integral morphism is integral so in order to prove (2)
it suffices to show that an integral morphism is closed.
This follows from Algebra, Lemmas \ref{algebra-lemma-integral-going-up}
and \ref{algebra-lemma-going-up-closed}.

\medskip\noindent
Assume (2). We may assume $f$ is the morphism
$f : \Spec(A) \to \Spec(R)$ coming from a ring map
$R \to A$. Let $a$ be an element of $A$. We have to show
that $a$ is integral over $R$, i.e. that in the kernel
$I$ of the map $R[x] \to A$ sending $x$ to $a$ there is
a monic polynomial. Consider the ring
$B = A[x]/(ax -1)$ and let $J$ be the kernel of the
composition $R[x]\to A[x] \to B$. If $f\in J$ there exists
$q\in A[x]$ such that $f = (ax-1)q$ in $A[x]$ so if
$f = \sum_i f_ix^i$ and $q = \sum_iq_ix^i$, for all $i \geq 0$
we have $f_i = aq_{i-1} - q_i$. For $n \geq \deg q + 1$ the polynomial
$$
\sum\nolimits_{i \geq 0} f_i x^{n - i} =
\sum\nolimits_{i \geq 0} (a q_{i - 1} - q_i) x^{n - i} =
(a - x) \sum\nolimits_{i \geq 0} q_i x^{n - i - 1}
$$
is clearly in $I$; if $f_0 = 1$ this polynomial is also monic, so we
are reduced to prove that $J$ contains a polynomial with constant term $1$.
We do it by proving $\Spec(R[x]/(J + (x))$ is empty.


\medskip\noindent
Since $f$ is universally closed the base change $\Spec(A[x]) \to \Spec(R[x])$
is closed. Hence the image of the closed subset $\Spec(B) \subset \Spec(A[x])$
is the closed subset $\Spec(R[x]/J) \subset \Spec(R[x])$, see
Example \ref{example-scheme-theoretic-image} and
Lemma \ref{lemma-quasi-compact-scheme-theoretic-image}.
In particular $\Spec(B) \to \Spec(R[x]/J)$ is surjective. Consider the
following diagram where every square is a pullback:
$$
\xymatrix{
\Spec(B) \ar@{->>}[r]^g &
\Spec(R[x]/J)  \ar[r] &
\Spec(R[x])\\
\emptyset \ar[u] \ar[r] &
\Spec(R[x]/(J + (x)))\ar[u] \ar[r] &
\Spec(R) \ar[u]^0
}
$$
The bottom left corner is empty because it is the spectrum of
$R\otimes_{R[x]} B$ where the map $R[x]\to B$ sends $x$ to an
invertible element and $R[x]\to R$ sends $x$ to $0$. Since $g$
is surjective this implies $\Spec(R[x]/(J + (x)))$ is empty, as
we wanted to show.
\end{proof}

\begin{lemma}
\label{lemma-integral-fibres}
Let $f : X \to S$ be an integral morphism.
Then every point of $X$ is closed in its fibre.
\end{lemma}

\begin{proof}
See Algebra, Lemma \ref{algebra-lemma-integral-no-inclusion}.
\end{proof}

\begin{lemma}
\label{lemma-integral-dimension}
Let $f : X \to Y$ be an integral morphism. Then $\dim(X) \leq \dim(Y)$.
If $f$ is surjective then $\dim(X) = \dim(Y)$.
\end{lemma}

\begin{proof}
Since the dimension of $X$ and $Y$ is the supremum of the dimensions
of the members of an affine open covering, we may assume $Y$ and $X$
are affine. The inequality follows from
Algebra, Lemma \ref{algebra-lemma-integral-dim-up}.
The equality then follows from
Algebra, Lemmas \ref{algebra-lemma-dimension-going-up}
and \ref{algebra-lemma-integral-going-up}.
\end{proof}

\begin{lemma}
\label{lemma-finite-quasi-finite}
A finite morphism is quasi-finite.
\end{lemma}

\begin{proof}
This is implied by Algebra, Lemma \ref{algebra-lemma-quasi-finite}
and Lemma \ref{lemma-quasi-finite-locally-quasi-compact}.
Alternatively, all points in fibres are closed points by
Lemma \ref{lemma-integral-fibres} (and the fact that a finite
morphism is integral) and use
Lemma \ref{lemma-quasi-finite-at-point-characterize} (3) to
see that $f$ is quasi-finite at $x$ for all $x \in X$.
\end{proof}

\begin{lemma}
\label{lemma-finite-proper}
Let $f : X \to S$ be a morphism of schemes. The following are equivalent
\begin{enumerate}
\item $f$ is finite, and
\item $f$ is affine and proper.
\end{enumerate}
\end{lemma}

\begin{proof}
This follows formally from
Lemma \ref{lemma-integral-universally-closed},
the fact that a finite morphism is integral and separated,
the fact that a proper morphism is the same thing as
a finite type, separated, universally closed morphism,
and the fact that an integral morphism of finite type is
finite (Lemma \ref{lemma-finite-integral}).
\end{proof}

\begin{lemma}
\label{lemma-closed-immersion-finite}
A closed immersion is finite (and a fortiori integral).
\end{lemma}

\begin{proof}
True because a closed immersion is affine
(Lemma \ref{lemma-closed-immersion-affine})
and a surjective ring map is finite and integral.
\end{proof}

\begin{lemma}
\label{lemma-finite-union-finite}
Let $X_i \to Y$, $i = 1, \ldots, n$ be finite morphisms of schemes.
Then $X_1 \amalg \ldots \amalg X_n \to Y$ is finite too.
\end{lemma}

\begin{proof}
Follows from the algebra fact that if $R \to A_i$, $i = 1, \ldots, n$
are finite ring maps, then $R \to A_1 \times \ldots \times A_n$ is finite too.
\end{proof}

\begin{lemma}
\label{lemma-finite-permanence}
Let $f : X \to Y$ and $g : Y \to Z$ be morphisms.
\begin{enumerate}
\item If $g \circ f$ is finite and $g$ separated then $f$ is finite.
\item If $g \circ f$ is integral and $g$ separated then $f$ is integral.
\end{enumerate}
\end{lemma}

\begin{proof}
Assume $g \circ f$ is finite (resp.\ integral) and $g$ separated.
The base change $X \times_Z Y \to Y$ is finite (resp.\ integral) by
Lemma \ref{lemma-base-change-finite}.
The morphism $X \to X \times_Z Y$ is
a closed immersion as $Y \to Z$ is separated, see
Schemes, Lemma \ref{schemes-lemma-section-immersion}.
A closed immersion is finite (resp.\ integral),
see Lemma \ref{lemma-closed-immersion-finite}.
The composition of finite (resp.\ integral) morphisms is finite
(resp.\ integral),
see Lemma \ref{lemma-composition-finite}. Thus we win.
\end{proof}

\begin{lemma}
\label{lemma-finite-monomorphism-closed}
Let $f : X \to Y$ be a morphism of schemes.
If $f$ is finite and a monomorphism, then $f$ is a closed immersion.
\end{lemma}

\begin{proof}
This reduces to
Algebra, Lemma \ref{algebra-lemma-finite-epimorphism-surjective}.
\end{proof}

\begin{lemma}
\label{lemma-finite-projective}
A finite morphism is projective.
\end{lemma}

\begin{proof}
Let $f : X \to S$ be a finite morphism. Then $f_*\mathcal{O}_X$ is
a quasi-coherent $\mathcal{O}_S$-module
(Lemma \ref{lemma-affine-equivalence-algebras})
of finite type
(by our definition of finite morphisms and
Properties, Lemma \ref{properties-lemma-finite-type-module}).
We claim there is a closed immersion
$$
\sigma :
X
\longrightarrow
\mathbf{P}(f_*\mathcal{O}_X) =
\underline{\text{Proj}}_S(\text{Sym}^*_{\mathcal{O}_S}(f_*\mathcal{O}_X))
$$
over $S$, which finishes
the proof. Namely, we let $\sigma$ be the morphism which corresponds
(via Constructions, Lemma \ref{constructions-lemma-apply-relative})
to the surjection
$$
f^*f_*\mathcal{O}_X \longrightarrow \mathcal{O}_X
$$
coming from the adjunction map $f^*f_* \to \text{id}$. Then $\sigma$
is a closed immersion. Namely, affine locally on $S$ we can write
$X = \Spec(A)$ and $S = \Spec(R)$. Since $X$ is finite over $S$
we may choose $a_1, \ldots, a_n \in A$ generating $A$ as an $R$-module.
Then the $R$-algebra map $R[T_0, T_1, \ldots, T_n] \to \text{Sym}_R^*(A)$
sending $T_0$ to $1$ and $T_i$ to $a_i$ for $1 \leq i \leq n$ is surjective.
Whence $\mathbf{P}(f_*\mathcal{O}_X)$ is a closed subscheme of
$\mathbf{P}^n_S$ by Constructions, Lemma
\ref{constructions-lemma-surjective-graded-rings-map-proj}.
Thus it suffices to prove that the induced morphism
$X \to \mathbf{P}^n_S$ is a closed immersion
(see for example, Lemma \ref{lemma-closed-immersion-ideals}).
The reader checks that $X \to \mathbf{P}^n_S$ has image contained
in the open $D_+(T_0) \cong \mathbf{A}^n_S$ and that $X \to \mathbf{A}^n_S$
corresponds to the surjective $R$-algebra map
$R[x_1, \ldots, x_n] \to \text{Sym}_R^*(A)$ sending $x_i$ to $a_i$.
Whence $X \to \mathbf{P}^n_S$ is an immersion (as a composition
of a closed immersion and an open immersion). Since
$X$ is finite over $S$ the image of this immersion is closed
(this uses Lemma \ref{lemma-finite-proper},
Lemma \ref{lemma-image-proper-scheme-closed}, and
Constructions, Lemma \ref{constructions-lemma-projective-space-separated}).
We conclude by Schemes, Lemma \ref{schemes-lemma-immersion-when-closed}.
\end{proof}







